% ECONOMICS_PROBLEM: {"id": "G14-255", "title": "Passive investing and price discovery", "jel_code": "G14", "source_id": "OP-0315"}
% !TeX program = xelatex
\documentclass[11pt,a4paper]{article}
\usepackage[margin=23mm]{geometry}
\usepackage{fontspec}
\setmainfont{Latin Modern Roman}
\usepackage{amsmath,amssymb,mathtools}
\usepackage{xcolor,enumitem,booktabs,longtable,array,xurl,hyperref}
\definecolor{muted}{HTML}{53636C}
\hypersetup{colorlinks=true,urlcolor=blue,linkcolor=blue}
\setlength{\parindent}{0pt}
\setlength{\parskip}{5pt}
\newcommand{\R}{\mathbb R}
\newcommand{\E}{\mathbb E}
\newcommand{\N}{\mathbb N}
\newcommand{\ind}{\mathbf 1}
\newcommand{\Var}{\operatorname{Var}}
\newcommand{\Cov}{\operatorname{Cov}}
\newcommand{\supp}{\operatorname{supp}}
\DeclareMathOperator*{\argmin}{arg\,min}
\DeclareMathOperator*{\argmax}{arg\,max}
\newcommand{\entrymeta}[3]{{\sffamily\small\color{muted}\textbf{\texttt{#1}}\quad|\quad #2\par #3\par}}
\newcommand{\Problemref}[1]{\texttt{#1}}
\newcommand{\JELref}[2]{\texttt{#2} (JEL \texttt{#1})}
\begin{document}
\section*{Passive investing and price discovery}
\textbf{Problem ID:} \texttt{G14-255}\quad \textbf{JEL:} \texttt{G14}\par
% BEGIN ECONOMICS STATEMENT
\label{problem:OP-0315}
{\sffamily\small\color{muted}Original subject group: Asset pricing and financial markets\par}
\label{definitions:problem:30007597}
\textbf{Audit classification:} Background; proposed extension. The accessible review abstract discusses passive investing as a macrostructure example. It does not verify the exact informativeness/adoption question as open; that formulation is editorial.
\subsection*{Definitions and precise research target}
Fix a stock market with firms \(j=1,\ldots,J\), dividends \(D_{jt}\), prices \(P_{jt}\), active investors and index funds. Define the passive share \(s\in[0,1]\) as the fraction of invested equity wealth committed to a publicly specified index rule. Intervention \(do(s)\) changes this share while permitting active information acquisition, arbitrage capital and firm investment to adjust in equilibrium. Let \(V_{jt}=V_{jt}(s)\) be the discounted value of future distributions under a specified fundamental benchmark. Measure price informativeness by
\[I(s)=1-\frac{E[(V_{jt}-E[V_{jt}\mid P_{jt},\mathcal F_t^{\mathrm{public}}])^2]}{\operatorname{Var}(V_{jt})},\]
where \(\mathcal F_t^{\mathrm{public}}\) is public information, \(\operatorname{Var}(V_{jt})>0\), and all variances use a fixed population. The task is to estimate the causal response of \(I(s)\), valuation deviations and investment allocation to a sustained increase in passive ownership. Define valuation deviations using a separately disciplined benchmark, not prices themselves. Distinguish a reallocation of wealth from an index-composition change. Establish identifying assumptions for index inclusion or fund-flow variation and characterize spillovers to untreated firms. Quantify how endogenous active-investor entry alters effects at larger passive shares; local estimates need not extrapolate to complete passivity.

\textbf{Additional formal definitions and scope (editorial).} Fix a sampling law over firm-date pairs and square-integrable \(V_{jt}(s)\). A fundamental benchmark specifies discounting, information and terminal-value or transversality assumptions independently of the observed price. Conditional expectation in the informativeness statistic uses the equilibrium law under the intervention. Its variance denominator is \(\operatorname{Var}(V_{jt}(s))>0\), giving \(0\leq I(s)\leq1\); if a fixed reference variance is used instead, the interpretation and these bounds change. Public information already makes \(I(s)\) positive in general: incremental information from prices can instead be measured by the reduction in mean squared error relative to \(E[V\mid\mathcal F^{\rm public}]\). Index rules specify constituents, weights and rebalancing dates; passive wealth share and passive fraction of individual-stock ownership differ. Define active information costs and entry, and choose or bound the equilibrium when participation produces multiplicity.

For the identification part, declare an admissible class \(\Theta\) of structural models, including preferences, technologies, shock laws, measurement/sampling rules and any equilibrium selector. If \(O\) denotes the complete observable history and \(T(\theta)\) the target defined above, the identified set is \(\mathcal I_T=\{T(\theta):\theta\in\Theta,\ \mathcal L_\theta(O)=\mathcal L_{\rm obs}(O)\}\); point identification means this set is a singleton. A \(do(a)\) comparison replaces stated policy/technology equations while fixing the declared exogenous law and re-solving the remaining equations. These definitions do not assert that an identification theorem holds without further restrictions.
\subsection*{Variants and assumption changes}
\begin{enumerate}
\item \textbf{Fixed information production (editorial).} Change passive wealth share while holding active signal precision and capital fixed; estimate this deliberately partial-equilibrium effect.
\item \textbf{Endogenous information and entry (editorial).} Allow active investors to respond and firms to invest; characterize informativeness and allocation over sustained passive shares, including equilibrium multiplicity.
\end{enumerate}
\subsection*{References and source audit}
\begin{enumerate}
\item Publisher metadata and abstract verified. Locator: abstract. Background only; no exact open passage verified. URL: \url{https://www.annualreviews.org/content/journals/10.1146/annurev-financial-090524-120754}.
\end{enumerate}
\begin{enumerate}\raggedright
\item\hypertarget{definitions:source:30007597:1}{} Valentin Haddad; Tyler Muir. 2025. \emph{Market Macrostructure: Institutions and Asset Prices}. Publisher abstract.
\url{https://www.annualreviews.org/content/journals/10.1146/annurev-financial-090524-120754}.
DOI: \url{https://doi.org/10.1146/annurev-financial-090524-120754}.
\end{enumerate}

\subsection*{Common conventions: Source mathematical conventions}

These conventions apply to each entry unless explicitly replaced.

\textbf{Models and identification.} A primitive model \(m\) specifies all probability laws, feasible choices, information, equilibrium and measurement rules used by its target. The admissible class \(\mathcal M\) is fixed independently of outcomes. If \(P_O(m)\) is its observable probability law and \(\tau(m)\) the target, its identified set at observable law \(P\) is
\[
 \mathcal I_\tau(P)=\{\tau(m):m\in\mathcal M,\ P_O(m)=P\}.
\]
Point identification means that this set is a singleton. An empty set signals incompatibility of data and assumptions. Coordinate bounds need not describe the joint set. Bounds are sharp only if every retained target is attainable or arbitrarily approachable by compatible models. Finite bounds require explicit restrictions on unobserved quantities; unrestricted confounding, hidden wealth, extrapolation or arbitrary equilibrium selection can yield uninformative sets.

\textbf{Causal interventions.} A potential outcome \(Y_i(p)\) is the outcome under a complete policy regime \(p\), including timing, financing, information and market spillovers. The target population and units are fixed before treatment. With interference, use the complete assignment vector or a justified exposure mapping; individual no-interference notation alone cannot identify an economy-wide intervention. A difference of mean potential outcomes does not require the cross-world joint distribution. Conditioning on post-treatment migration, survival, enrollment or employment changes the estimand and needs separate justification.

\textbf{Statistical statements.} An estimator, confidence set, forecast or test specifies a sampling law, model class and loss/coverage criterion. Uniform \((1-\alpha)\)-coverage means \(\inf_{m\in\mathcal M}\Pr_m\{\tau(m)\in C_n\}\geq1-\alpha\), or an explicitly asymptotic version. The whole parameter space trivially has coverage, so informativeness requires a stated width, risk or power objective. A forecasting improvement is relative to a named benchmark, information set and evaluation population.

\textbf{Equilibrium and optimization.} An equilibrium comprises feasible plans, optimality, beliefs where needed, budget constraints and all market/resource clearing. The primitives or a stated benchmark define each correspondence \(\mathcal E(p;m)\). Nonemptiness is assumed or proved, never inferred just from compact policy sets. A selected-equilibrium target uses a declared selection rule; a robust target quantifies over all permitted equilibria. Use suprema or infima unless attainment is established. An \(\varepsilon\)-optimal policy is feasible and within \(\varepsilon\) of the finite optimum value. Report an empty feasible set separately.

\textbf{Welfare and accounting.} Normative utility, interpersonal normalization, social weights, numeraire, discounting and the use of public receipts are primitives. Behavioral utility need not coincide with the welfare criterion. Household welfare includes ownership income and rents once; adding profits again to compensating variation generally double-counts them. A monetary equivalent is defined by an explicit transfer and price convention, with monotonicity and range conditions. Average effects, mechanism contributions, transfers and welfare losses are different targets. Infinite sums require convergence and dynamic feasible plans require terminal or no-Ponzi conditions.

\textbf{Source versions.} A working-paper date, revision date and journal publication year are distinguished. A DOI for a journal version is not automatically the DOI of an earlier manuscript. Locators refer to the stated version and printed pages where specified. A metadata-only check does not verify a claim about a full-text conclusion. Every entry's source subsection narrows the provenance claim accordingly.
% END ECONOMICS STATEMENT
\end{document}
